\documentclass[11pt]{article}

\usepackage[margin=1in]{geometry}
\usepackage{amsmath,amssymb,amsthm}
\usepackage{enumitem}

\newtheorem{corollary}{Corollary}
\newtheorem{example}{Example}

\newcommand{\LRF}{\mathrm{LRF}}
\newcommand{\OPT}{\mathrm{OPT}}
\newcommand{\Cmax}{C_{\max}}
\newcommand{\Cmin}{C_{\max}^{\min}}

\begin{document}

\title{An Alternative Proof of the Kawaguchi--Kyan Bound for the
Largest-Ratio-First Rule}
\author{Uwe Schwiegelshohn\\
\small Robotics Research Institute, TU Dortmund University, D-44221 Dortmund, Germany}
\date{}
\maketitle

\begin{abstract}
In a well-known paper from 1986, Kawaguchi and Kyan proved that the
Largest-Ratio-First rule for total weighted completion time scheduling on
identical parallel machines has a worst case ratio of
\[
        \frac12(1+\sqrt{2}).
\]
We provide an alternative and simple proof for this result.
\end{abstract}

\noindent\textbf{Keywords:}
Total weighted completion time scheduling; worst case analysis; performance
guarantee.

\bigskip

An instance of the scheduling problem
\[
        P\|\sum_j w_j C_j
\]
consists of $n$ jobs. Each job has a weight $w_j \in \mathbb{R}^+$ and a
processing time $p_j \in \mathbb{R}^+$. The goal is to schedule the jobs on
$m$ identical parallel machines such that the weighted sum of the resulting
job completion times $C_j$ (total weighted completion time) is minimized.
List scheduling is frequently applied to generate schedules for this type of
problem: at time 0 or whenever a machine becomes idle, the available jobs are
scanned in a previously established order to determine the first job that can
be started on machines, idle at this time. Any list schedule for the
$P\|\sum_j w_j C_j$ problem has no intermediate idle times as all jobs become
available at time 0 and each job only requires a single machine. Delaying the
execution of a job while a machine is idle increases the completion time of
this job while it does not reduce the completion time of any other job.
Hence, for each instance of the $P\|\sum_j w_j C_j$ problem, there is a job
order such that the corresponding list schedule is an optimal schedule. An
LRF (Largest-Ratio-First) schedule is a specific list schedule that uses a
decreasing order of the ratio $w_j/p_j$ for the jobs. In a celebrated paper
from 1986, Kawaguchi and Kyan~[1] showed that the total weighted completion
time of an LRF schedule is at most a factor
\[
        \frac12(1+\sqrt{2})
\]
larger than the optimal total weighted completion time; furthermore this
bound is tight.

Unfortunately, the original proof by Kawaguchi and Kyan~[1] is rather long
and involved. Therefore, this note provides an alternative and simpler
version of this proof divided into four corollaries. Step by step,
Corollaries 1, 3 and 5 simplify the performance analysis by reducing the set
of instances that must be examined. First, Corollary 1 shows that the worst
case performance occurs in an instance in which $w_j=p_j$ holds for all jobs.
For these instances, Queyranne~[2] has given a simple expression to calculate
the total weighted completion time in a single machine schedule without
intermediate idle times. For the sake of completeness, we state this result
in Corollary 2 and give a brief proof before extending the expression to
parallel identical machines. For the purpose of our performance analysis, we
show in Corollary 3 that it is sufficient to consider instances that contain
at most $m-1$ so-called long jobs while the processing time of all other jobs
is arbitrarily small. Finally, Corollary 5 enables us to additionally demand
that all long jobs have the same processing time. As a result of these four
corollaries, we can express the worst case performance with two functions of
two variables each. Determining the maximum of these functions within the
permitted ranges of the variables yields the performance guarantee of an LRF
schedule.

Before starting with the actual proof, we briefly point out the convexity of
the square function as this property is repeatedly used in the rest of this
note:
\[
        (x-z)^2+(y+z)^2
        =
        x^2+y^2-2z(x-y-z)
        <
        x^2+y^2
\]
for three positive numbers $x$, $y$, and $z$ with $x-y>z$. Therefore, for
$k$ numbers $x_i$ with a fixed sum $\sum_{i=1}^k x_i$, the sum of squares
$\sum_{i=1}^k x_i^2$ becomes minimal if
\[
        x_i=\frac1k\sum_{i=1}^k x_i
\]
holds for all $1\leq i\leq k$.

We extend our previous definition of an instance $I$ for the
$P\|\sum_j w_j C_j$ problem to additionally include an LRF order for all
jobs. This LRF order produces the so-called primary LRF schedule
$S_{\LRF}(I)$. Therefore, every instance has a unique primary LRF schedule
$S_{\LRF}$ even if all jobs have the same $w_j/p_j$ ratio. Throughout the
proof, we say that $C(I)$ is the total weighted completion time of
$S_{\LRF}(I)$ while $C^*(I)$ is the total weighted completion time of an
optimal schedule $S_{\OPT}(I)$ for instance $I$. Furthermore, we use the
notation
\[
        \lambda(I)=\frac{C(I)}{C^*(I)}
\]
to describe the performance ratio between $S_{\LRF}(I)$ and $S_{\OPT}(I)$.

\begin{corollary}
For every instance $I$ of the $P\|\sum_j w_j C_j$ problem, there is an
instance $I'$ of the problem $P\|\sum_j p_j C_j$ with
$\lambda(I')\geq \lambda(I)$.
\end{corollary}

\begin{proof}
Let $r_1>r_2>\cdots>r_d$ be the $d$ different $w_j/p_j$ ratios of an
instance $I$ and $r_{d+1}=0$. Note that we can use the term
\[
        r_i=\sum_{k=i}^d (r_k-r_{k+1}).
\]
Furthermore, let $i(j)$ denote the index of the ratio for job $j$, that is
$w_j/p_j=r_{i(j)}$. Then we write
\[
        w_j
        =
        r_{i(j)}p_j
        =
        \left(\sum_{k=i(j)}^d (r_k-r_{k+1})\right)p_j
\]
and use this term to express the total weighted completion time of an
arbitrary schedule for instance $I$:
\[
\begin{aligned}
        \sum_{j\in I} w_j C_j
        &=
        \sum_{j\in I}
        \left(
            \sum_{k=i(j)}^d (r_k-r_{k+1})
        \right)p_j C_j                                     \\
        &=
        \sum_{k=1}^d (r_k-r_{k+1})
        \sum_{j\in I\mid i(j)\leq k} p_j C_j .
\end{aligned}
\]
Next, we define a sequence of instances
\[
        I_k=\{j\in I\mid i(j)\leq k\}
\]
for $1\leq k\leq d$ and set $w_j=p_j$ for every job in any $I_k$. Note that
a job belongs to $I_k$ if it precedes a job $j\in I_k$ in any LRF order of
instance $I$. For the jobs in any $I_k$, we use the LRF order of these jobs
in instance $I$, as our weight selection allows every order of the jobs in
$I_k$. Therefore, the completion times of any job $j\in I_k$ in
$S_{\LRF}(I_k)$ and in $S_{\LRF}(I)$ are identical. With $C_j^*$ denoting the
completion time of job $j$ in an optimal schedule for instance $I$, we have
\[
\begin{aligned}
        \lambda(I)
        &=
        \frac{\sum_{k=1}^d (r_k-r_{k+1})C(I_k)}
        {\sum_{k=1}^d
        \left((r_k-r_{k+1})\sum_{j\in I_k}p_jC_j^*\right)}
        \\
        &\leq
        \frac{\sum_{k=1}^d (r_k-r_{k+1})C(I_k)}
        {\sum_{k=1}^d (r_k-r_{k+1})C^*(I_k)}
        \leq
        \max_k \lambda(I_k).
        \qedhere
\end{aligned}
\]
\end{proof}

For single machine list schedules of instances in which $w_j=p_j$ holds for
all jobs, Queyranne~[2] has presented a simple formula to determine the total
weighted completion time with $\Cmax(S)$, denoting the maximum completion
time or makespan of schedule $S$:

\begin{corollary}
For any schedule $S$ on a single machine without intermediate idle times,
\[
        \sum_{j\in I}p_j C_j
        =
        \frac12\left(\sum_{j\in I}p_j^2+(\Cmax(S))^2\right)
\]
holds.
\end{corollary}

\begin{proof}
Remember the relation
\[
\begin{aligned}
        \left(\sum_{i=1}^n a_i\right)^2
        &=
        \sum_{i=1}^n a_i\sum_{k=1}^n a_k                         \\
        &=
        \sum_{i=1}^n a_i^2
        +2\sum_{i=2}^n a_i\sum_{k=1}^{i-1}a_k                    \\
        &=
        2\sum_{i=1}^n a_i\sum_{k=1}^{i}a_k-\sum_{i=1}^n a_i^2 .
\end{aligned}
\]
Assume that in schedule $S$, the jobs in instance $I$ are scheduled in order
of increasing index. Then, we have
\[
\begin{aligned}
        \sum_{j=1}^n p_j C_j
        &=
        \sum_{j=1}^n p_j\sum_{k=1}^j p_k                         \\
        &=
        \frac12\sum_{j=1}^n p_j^2
        +\frac12\left(\sum_{j=1}^n p_j\right)^2                  \\
        &=
        \frac12\left(\sum_{j\in I}p_j^2+(\Cmax(S))^2\right).
        \qedhere
\end{aligned}
\]
\end{proof}

Using the notation $C_{\max}^i(S)$ to describe the last completion time or
makespan of a list schedule $S$ on machine $i$ in the parallel identical
machine scenario, we simply apply Corollary 2 to every machine separately and
combine the squares of the processing times into a single term:
\begin{equation}
        \sum_{j\in I} p_j C_j
        =
        \frac12
        \left(
            \sum_{j\in I}p_j^2
            +
            \sum_{i=1}^m (C_{\max}^i(S))^2
        \right).
\end{equation}
Furthermore, we introduce
\[
        \Cmin(S)=\min_i\{C_{\max}^i(S)\}
\]
to be the minimum value of the machine makespans in the parallel identical
machine scenario. Note that in any list schedule $S$, no job can start and at
most $m-1$ jobs can complete after $\Cmin(S)$.

\begin{corollary}
For every instance $I$ of the $P\|\sum_j p_j C_j$ problem, there is an
instance $I'$ of the $P\|\sum_j p_j C_j$ problem such that
$\lambda(I')\geq\lambda(I)$ holds and every job in instance $I'$ has an
arbitrarily small processing time if it starts before
$\Cmin(S_{\LRF}(I'))$.
\end{corollary}

\begin{proof}
Assume an arbitrary list schedule $S(I)$ for an instance $I$ of the
$P\|\sum_j p_j C_j$ problem. We generate an instance $J$ by taking a job
$j_0$ from $I$ and replacing it in schedule $S(I)$ with two jobs $j_1$ and
$j_2$ such that $p_{j_1}+p_{j_2}=p_{j_0}$, $C_{j_2}=C_{j_0}$, and
$C_{j_1}=C_{j_0}-p_{j_2}$ hold. Note that this exchange does not alter the
completion time of any other job in schedule $S(I)$. Therefore, we have
\[
        C_{\max}^i(S(I))=C_{\max}^i(S(J))
\]
for all $1\leq i\leq m$. Then Eq.~(1) produces
\[
        \sum_{j\in I}p_j C_j-\sum_{j\in J}p_j C_j
        =
        \frac12(p_{j_0}^2-p_{j_1}^2-p_{j_2}^2)
        =
        p_{j_1}p_{j_2}.
\]
If $C_{j_1}\leq \Cmin(S_{\LRF}(I))$ holds then the resulting schedule
$S_{\LRF}(J)$ for instance $J$ is a list schedule. As any job order for
instance $J$ is an LRF order, we can select an order for instance $J$ such
that we have $S(J)=S_{\LRF}(J)$. As $S_{\OPT}(J)$ is also a list schedule, we
have
\[
        1\leq \lambda(I)
        =
        \frac{C(I)}{C^*(I)}
        \leq
        \frac{C(I)-p_{j_1}p_{j_2}}{C^*(I)-p_{j_1}p_{j_2}}
        \leq
        \frac{C(J)}{C^*(J)}
        =
        \lambda(J).
\]
We repeat this job splitting until all jobs starting before
$\Cmin(S_{\LRF}(I))$ have an arbitrarily small processing time. As at most
$m-1$ jobs can start at $\Cmin(S_{\LRF}(I))$, all but at most $m-1$ jobs have
an arbitrarily small processing time.
\end{proof}

Due to Corollaries 1 and 3, we consider from now on only instances $I$ in
which $w_j=p_j$ holds for all jobs and all jobs have an arbitrarily small
processing time with the exception of at most $m-1$ so-called long jobs that
start at $\Cmin(S_{\LRF}(I))$. We use $L(I)\subseteq I$ to denote the set of
these long jobs in instance $I$. Furthermore, we introduce the set
$K(I)\subseteq L(I)$ of long jobs to be the largest set such that
\[
        p_j\geq \frac1{m-|K(I)|}\sum_{j\in I\setminus K(I)}p_j
\]
holds for all $j\in K(I)$. We obtain $K(I)$ by adding long jobs in decreasing
order of their processing times and testing the property of $K(I)$. All jobs
from $I\setminus L(I)$ are called small jobs. Note that an instance $I$ only
contains small jobs if $\Cmin(S_{\LRF}(I))>0$ holds.

As the processing time of each small job is arbitrarily small, we can ignore
the term $\sum_{j\in I\setminus L(I)}p_j^2$ in Eq.~(1) for any list schedule
$S(I)$ leading to
\begin{equation}
        \sum_{j\in I}p_j C_j
        =
        \frac12
        \left(
            \sum_{j\in L(I)}p_j^2
            +
            \sum_{i=1}^m(C_{\max}^i(S(I)))^2
        \right)
\end{equation}
and
\begin{equation}
\begin{aligned}
        C(I)
        &=
        \frac12\sum_{j\in L(I)}p_j^2
        +
        \frac12\sum_{j\in L(I)}
            (p_j+\Cmin(S_{\LRF}(I)))^2                    \\
        &\quad
        +
        \frac12\left((m-|L(I)|)(\Cmin(S_{\LRF}(I)))^2\right).
\end{aligned}
\end{equation}
Note that Eq.~(3) is the application of Eq.~(2) to $S_{\LRF}(I)$.

Next, we address the generation of an optimal schedule for our instances.
Consider a job order for instance $I$ in which the long jobs precede all small
jobs. In the corresponding list schedule $S_{\OPT}(I)$, all long jobs are
started at time 0 due to $|L(I)|<m$. Clearly, we have
\[
        \Cmin(S_{\OPT}(I))
        \leq
        \frac1{m-|K(I)|}\sum_{j\in I\setminus K(I)}p_j .
\]
If a small job $j_0$ finishes after
\[
        \frac1{m-|K(I)|}\sum_{j\in I\setminus K(I)}p_j ,
\]
we apply job splitting to $j_0$ at time
\[
        \frac1{m-|K(I)|}\sum_{j\in I\setminus K(I)}p_j
\]
as described in the proof of Corollary 3. In the list for schedule
$S_{\OPT}(I)$, the first new job $j_1$ replaces $j_0$ and the other new job
$j_2$ is put at the end of the list. After at most $m-|K(I)|-1$ such
splitting steps,
\[
        C_{\max}^i(S_{\OPT}(I))
        =
        \Cmin(S_{\OPT}(I))
        =
        \frac1{m-|K(I)|}\sum_{j\in I\setminus K(I)}p_j
\]
holds if machine $i$ executes a small job. Therefore, every machine $i$ with
$C_{\max}^i(S_{\OPT}(I))>\Cmin(S_{\OPT}(I))$ can only execute a single long
job. Due to Eq.~(2) and the convexity of the square function, it is not
possible to further reduce
\[
        \sum_{i=1}^m(C_{\max}^i(S_{\OPT}(I)))^2
\]
by transferring any jobs. Hence, schedule $S_{\OPT}(I)$ is an optimal
schedule for instance $I$.

Now, we introduce a transformation that replaces a set $M$ of positive
numbers with a number $p$ that is used $t$ times such that the following
equations hold:
\begin{align}
        \sum_{i\in M} i &= t\cdot p, \\
        \sum_{i\in M} i^2 &= t\cdot p^2 .
\end{align}
Eqs.~(4) and (5) lead to a unique solution with
\[
        t=\frac{\left(\sum_{i\in M}i\right)^2}{\sum_{i\in M}i^2}
        \qquad\text{and}\qquad
        p=\frac{\sum_{i\in M}i^2}{\sum_{i\in M}i}.
\]
In Eq.~(2), this transformation is either applied to the set $L(I)$ of long
jobs or to some makespans. However, $t$ need not be an integer but may be a
rational number with denominator $q$ resulting in an instance with a
fractional number of jobs. To obtain an instance with an integral number of
jobs, we simply replace each machine and each job of the original instance
$I$ with $q$ identical copies. In the LRF order, all replicas of a job occur
one after another. Therefore, the schedule of each machine is replicated for
schedule $S_{\LRF}$. Replicating the schedule of each machine for schedule
$S_{\OPT}$ also yields an optimal schedule for the same reasons discussed
before. Hence, $C(I)$ and $C^*(I)$ are simply multiplied by $q$, and
$\lambda(I)$ remains unchanged.

The application of the transformation to the set $L(I)$ or to machine
makespans of a list schedule $S(I)$ does not change
$\sum_{i=1}^m C_{\max}^i(S(I))$ due to Eq.~(4). Furthermore, due to the
convexity of the square function, we have
\begin{equation}
\begin{aligned}
        t
        &=
        \frac{\left(\sum_{j\in L(I)}p_j\right)^2}
        {\sum_{j\in L(I)}p_j^2}                                      \\
        &\leq
        \frac{\left(\sum_{j\in L(I)}p_j\right)^2}
        {|L(I)|\left(\frac{\sum_{j\in L(I)}p_j}{|L(I)|}\right)^2}       \\
        &=
        \frac{|L(I)|\left(\sum_{j\in L(I)}p_j\right)^2}
        {\left(\sum_{j\in L(I)}p_j\right)^2}
        =
        |L(I)|.
\end{aligned}
\end{equation}

\begin{example}
Consider a set $L(I)$ of six long jobs with processing times $5,5,4,3,2$,
and $1$, respectively. Our transformation yields five long jobs, all with
processing time $4$ producing
\[
        5\cdot 4=20=\sum_{j\in L(I)}p_j
        \qquad\text{and}\qquad
        5\cdot 4^2=80=\sum_{j\in L(I)}p_j^2 .
\]
\end{example}

\begin{corollary}
For every instance $I$ of the $P\|\sum_j p_j C_j$ problem, there is an
instance $I'$ of the $P\|\sum_j p_j C_j$ problem with
$\lambda(I')\geq\lambda(I)$ such that instance $I'$ contains at most $m-1$
jobs with the same processing time while all its other jobs have an
arbitrarily small processing time.
\end{corollary}

\begin{proof}
Due to Corollary 3, assume schedule $S_{\LRF}(I)$ for instance $I$ in which
only jobs with an arbitrarily small processing time start before
$\Cmin(S_{\LRF}(I))$.

We apply our transformation to the set $L(I)$ to generate an instance $I'$ in
which all $t$ long jobs have the same processing time. In the LRF order of
both instances, long jobs appear after all small jobs. Therefore, we have
\[
        \Cmin(S_{\LRF}(I))=\Cmin(S_{\LRF}(I')).
\]
We discuss the consequences of the transformation to the primary LRF schedule
and an optimal schedule separately. Due to Eqs.~(3)--(5), we have
\[
\begin{aligned}
        2C(I)
        &=
        \sum_{j\in L(I)}p_j^2
        +
        \sum_{j\in L(I)}(p_j+\Cmin(S_{\LRF}(I)))^2                  \\
        &\quad
        +(m-|L(I)|)(\Cmin(S_{\LRF}(I)))^2                           \\
        &=
        \sum_{j\in L(I)}p_j^2
        +
        \sum_{j\in L(I)}p_j^2
        +
        2\sum_{j\in L(I)}(p_j\cdot\Cmin(S_{\LRF}(I)))               \\
        &\quad
        +m(\Cmin(S_{\LRF}(I)))^2                                    \\
        &=
        t\cdot p^2+t\cdot p^2
        +2t\cdot p\cdot\Cmin(S_{\LRF}(I'))
        +m(\Cmin(S_{\LRF}(I')))^2                                  \\
        &=
        t\cdot p^2
        +t(p+\Cmin(S_{\LRF}(I')))^2
        +(m-t)(\Cmin(S_{\LRF}(I')))^2                              \\
        &=
        2C(I').
\end{aligned}
\]

Next we address the optimal schedules for instances $I$ and $I'$. To show
that $C^*(I')\leq C^*(I)$ holds, we consider two cases:

\begin{enumerate}[label=(\arabic*)]
\item
$C_{\max}^i(S_{\OPT}(I'))=\Cmin(S_{\OPT}(I'))$ holds for all machines $i$.

Remember that any list schedule $S(I')$ is optimal if
$C_{\max}^i(S(I'))=\Cmin(S(I'))$ holds for all machines $i$ due to the
convexity of the square function. Furthermore, the transformation does not
change $\sum_{i=1}^m C_{\max}^i(S(I))$ for any list schedule $S(I)$.
Therefore, Eqs.~(2) and (5) yield together with the convexity of the square
function:
\[
\begin{aligned}
        2C^*(I)
        &=
        \sum_{j\in L(I)}p_j^2
        +
        \sum_{i=1}^m(C_{\max}^i(S_{\OPT}(I)))^2                 \\
        &=
        t\cdot p^2
        +
        \sum_{i=1}^m(C_{\max}^i(S_{\OPT}(I)))^2                 \\
        &\geq
        t\cdot p^2
        +
        m(\Cmin(S_{\OPT}(I')))^2
        =
        2C^*(I').
\end{aligned}
\]

\item
$C_{\max}^i(S_{\OPT}(I'))>\Cmin(S_{\OPT}(I'))$ holds for some machines $i$.
In the following sequence of equations, we show that $C^*(I')\leq C^*(I)$
also holds for this case. The individual equations are later explained in
detail.
\begin{align}
        2C^*(I)
        &=
        \sum_{j\in L(I)}p_j^2
        +
        \sum_{j\in K(I)}p_j^2
        +
        (m-|K(I)|)(\Cmin(S_{\OPT}(I)))^2                         \tag{7}\\
        &=
        t\cdot p^2+t'\cdot(p')^2
        +(m-|K(I)|-k)(\Cmin(S_{\OPT}(I)))^2                      \tag{8}\\
        &\geq
        t\cdot p^2+t'\cdot(p')^2
        +(m-t')
        \left(
            \frac{m-|K(I)|-k}{m-t'}\Cmin(S_{\OPT}(I))
        \right)^2                                                \tag{9}\\
        &\geq
        t\cdot p^2+t\cdot p^2
        +(m-t)(\Cmin(S_{\OPT}(I')))^2
        =
        2C^*(I').                                                \tag{10}
\end{align}

First, we introduce the size
\[
        k=
        \frac{\sum_{j\in L(I)\setminus K(I)}p_j}
        {\Cmin(S_{\OPT}(I))}
        \leq
        |L(I)|-|K(I)|.
\]
If necessary, we use replication to guarantee an integral $k$. Eq.~(7) is
the application of Eq.~(2) on schedule $S_{\OPT}(I)$ by using set $K(I)$.
We apply our transformation twice to Eq.~(7):
\begin{enumerate}[label=(\alph*)]
\item on set $L(I)$ to transform $\sum_{j\in L(I)}p_j^2$;
\item on the makespan of all machines that execute jobs in $K(I)$ and of
$k$ additional machines with
$C_{\max}^i(S_{\OPT}(I))=\Cmin(S_{\OPT}(I))$.
\end{enumerate}
The second transformation produces the values $t'$ and $p'$. Remember that
we use replication to guarantee an integral $t'$. Due to the convexity of the
square function and $0\leq p_j\leq \Cmin(S_{\OPT}(I))$ for all jobs
$j\in L(I)\setminus K(I)$, we have
\begin{align}
        t'\cdot p'
        &=
        \sum_{j\in K(I)}p_j+k\cdot\Cmin(S_{\OPT}(I))
        =
        \sum_{j\in L(I)}p_j
        =
        t\cdot p,                                                \tag{11}\\
        t'\cdot(p')^2
        &=
        \sum_{j\in K(I)}p_j^2+k(\Cmin(S_{\OPT}(I)))^2             \notag\\
        &=
        \sum_{j\in K(I)}p_j^2+k(\Cmin(S_{\OPT}(I)))^2
        +(|L(I)|-|K(I)|-k)\cdot 0^2                              \notag\\
        &\geq
        \sum_{j\in K(I)}p_j^2
        +
        \sum_{j\in L(I)\setminus K(I)}p_j^2
        =
        \sum_{j\in L(I)}p_j^2
        =
        t\cdot p^2,                                              \tag{12}\\
        &\Rightarrow
        p'\geq p
        \quad\text{and}\quad
        t'\leq t.                                                \tag{13}
\end{align}

In Eq.~(8), the result of the second transformation to Eq.~(7) is shown.
Note that both transformations are formally applied and do not produce a
valid schedule. In particular, $|K(I)|+k-t'$ machines are not used in
Eq.~(8) as $t'\leq |K(I)|+k$ holds due to Eq.~(6).

In Eq.~(9), we evenly distribute the total load
\[
        (m-|K(I)|-k)\Cmin(S_{\OPT}(I))
\]
of all small jobs on $m-t'$ machines. This distribution reduces the total
weighted completion time due to the convexity of the square function.

Eq.~(10) shows schedule $S_{\OPT}(I')$. In this schedule, only
$m-t\leq m-t'$ machines are available to execute the small jobs due to
Eq.~(13). Therefore, we have
\[
        \frac{m-|K(I)|-k}{m-t'}\Cmin(S_{\OPT}(I))
        \leq
        C_{\max}^i(S_{\OPT}(I')).
\]
Due to this relation, Eqs.~(11)--(13), and the convexity of the square
function, the total weighted completion time is reduced again.
\end{enumerate}
\end{proof}

Due to Corollary 5, we can determine $\max_I\lambda(I)$ by examining only
instances and schedules with the following properties:
\begin{enumerate}[label=(\arabic*)]
\item An instance contains $u<m$ long jobs with processing time $p$ and
additional jobs with an arbitrary small processing time and a total workload
$m$.
\item For all jobs, $w_j=p_j$ holds.
\item In schedule $S_{\LRF}$, all small jobs are executed in the interval
$[0,1]$ while all long jobs are started at 1.
\item In schedule $S_{\OPT}$, all long jobs are started at 0 and
$C_{\max}^i(S_{\OPT})=\Cmin(S_{\OPT})$ holds for all machines that execute at
least one small job.
\end{enumerate}
As every list schedule without small jobs is always optimal, such instances
can be ignored. As we can divide the processing time of each job by a
constant without changing $\lambda(I)$, we define the total workload of all
small jobs to be $m$. As a result, we obtain a numerical optimization problem
with two variables $p>0$ and $0<u/m<1$ for each case of the proof of
Corollary 5. Note that Case 1 applies if and only if
\[
        m+u\cdot p\geq m\cdot p
        \quad\Longleftrightarrow\quad
        p\leq \frac1{1-u/m}.
\]
We solve the problems by determining the roots of the derivatives.

\[
\text{(1)}\qquad p\leq \frac1{1-u/m}.
\]
\[
\begin{aligned}
        \lambda(u,p)
        &=
        \frac{\frac12(u\cdot p^2+u\cdot(p+1)^2+m-u)}
        {\frac12\left(u\cdot p^2
        +m\left(\frac{u\cdot p+m}{m}\right)^2\right)}             \\
        &=
        \frac{2\frac{u}{m}p^2+2\frac{u}{m}p+1}
        {\frac{u}{m}p^2+\left(\frac{u}{m}\right)^2p^2
        +2\frac{u}{m}p+1}
        =
        f\left(\frac{u}{m},p\right).
\end{aligned}
\]
As there is no common root of both derivatives
$\partial f/\partial(u/m)$ and $\partial f/\partial p$ within the feasible
area of the function, we examine the roots of the derivative on the
boundaries
\[
        \left\{
            p=0,\ \frac{u}{m}=0,\ \frac{u}{m}=1,\
            p=\frac1{1-u/m}
        \right\}
\]
of the feasible area. As we have $f(u/m,p)=1$ for
$p=0$, $u/m=0$, and $u/m=1$, we must only consider
\[
        \frac{d f\left(u/m,\frac1{1-u/m}\right)}{d(u/m)}.
\]
This way, we obtain the maximum of
\[
        f\left(\frac{u}{m},p\right)=4-2\sqrt{2}
\]
for
\[
        \frac{u}{m}=\sqrt{2}-1
        \qquad\text{and}\qquad
        p=1+\frac{\sqrt{2}}2=\frac1{1-u/m}.
\]

\[
\text{(2)}\qquad p>\frac1{1-u/m}.
\]
\[
\begin{aligned}
        \lambda(u,p)
        &=
        \frac{\frac12(u\cdot p^2+u\cdot(p+1)^2+m-u)}
        {\frac12\left(2u\cdot p^2
        +(m-u)\left(\frac{m}{m-u}\right)^2\right)}                 \\
        &=
        \frac{2\frac{u}{m}p^2+2\frac{u}{m}p+1}
        {2\frac{u}{m}p^2+\frac1{1-u/m}}
        =
        f\left(\frac{u}{m},p\right).
\end{aligned}
\]
Note that the value on the boundary $p=1/(1-u/m)$ is the same as in the
first case while $f(u/m,p)=1$ holds for $p\to\infty$ and $u/m=0$.

There is a single common root of the derivatives
$\partial f/\partial(u/m)$ and $\partial f/\partial p$ at
\[
        \frac{u}{m}=1-\frac{\sqrt{2}}2
        \qquad\text{and}\qquad
        p=1+\sqrt{2}
        >
        \frac1{1-u/m}
        =
        \sqrt{2}.
\]
At this position, we obtain the maximum of
\[
        f\left(\frac{u}{m},p\right)
        =
        \frac12(1+\sqrt{2})
        >
        4-2\sqrt{2}.
\]
Finally note that the machine value $u/m=1-\sqrt{2}/2$ is not feasible as it
is an irrational number. But we can get arbitrarily close to it by using a
sufficiently large number of replications.

\section*{Acknowledgements}

The author thanks the referees for their valuable comments and suggestions
that helped to improve this note significantly.

\begin{thebibliography}{9}

\bibitem{KawaguchiKyan1986}
T. Kawaguchi, S. Kyan,
Worst case bound of an LRF schedule for the mean weighted flow-time problem,
\emph{SIAM Journal on Computing} 15 (4) (1986) 1119--1129.

\bibitem{Queyranne1993}
M. Queyranne,
Structure of a simple scheduling polyhedron,
\emph{Mathematical Programming} 58 (2) (1993) 263--285.

\end{thebibliography}

\end{document}
